\documentclass{article}
\usepackage{amsmath} % for mathematical symbols
\usepackage{graphicx} % for including images

\begin{document}
	
	\title{Answers \& Manual}
	\author{Paul Tr\"oger}
	
	\maketitle
	
	\section{Answers}
		\subsection{Task 1}
			\textbf{Which columns of adult\_small are identifiers, which are quasi-identifiers, and which are sensitive attributes?}\\
			There are no identifiers left in adult\_small because we removed them (Name, DOB and SSN).
			Everything that's neither sensitive nor an identifier is a quasi identifier, here: Education and Marital Status.
			The sensitive attribute is the Target column.
			
			\textbf{For which $k \in [1, 10]$ does adult\_small satisfy k-anonymity? For the cases that do not satisfy k-anonymity, why not?}\\
			Only $k$ = 1 is satisfied in adult\_small because we have many quasi-identifier combinations that hold only one person.
			For example there is only one person that has a Doctorate and was never married.
		\subsection{Task 2}
			\textbf{For which rows is a homogeneity attack possible, and why?}\\
			We can observe that every entry with  '<HS' is in target group '<=50k'.
			Because we choose to split target into only two groups and because the cutoff is relatively high we can easily know the target value if we know that a person is '<HS'. Of course because we only have two quasi identifiers it is also very easy to identify which target group a person belongs on all entries.
		\subsection{Task 4}
			\textbf{How many rows would we have to suppress in addition to the generalization to achieve k = 3 and k = 7?}\\
			See output file.
			\textbf{How many digits of Zip and Age can/should you generalize, and how does this affect the number of suppressed rows?}\\
			As a general rule, the more digits are generalized, the fewer entries need to be suppressed.
			If both digits of Age are generalized we loose a lot of information so we should try to keep that at 1.
			For the zip we can see that we move from the thousands to only hundreds of entries to suppress if we move from 2 to 3 digits to suppress.
			So we choose zip 3 digits and age 1 digit.  
		\subsection{Task 6}
			\textbf{Is the adult dataset $l$-diverse?}\\
			The dataset is not $l$ = 2 diverse. See output.
			
			\textbf{If not, which quasi-identifier groups are preventing $l$-diversity?}\\
			From the output we can see that the dataset becomes 2-diverse if we remove any of the three quasi-identifiers.
			
			\textbf{What is the difference between probabilistic and entropy $l$-diversity, if any?}\\
			On this dataset with these quasi-identifiers there is no difference.
		\subsection{Task 7}
			\textbf{What is the largest $l$ for the generalized dataset for probabilistic $l$-diversity and entropy
				$l$-diversity?}\\
				For both values it's ~1.1. See output.
			\textbf{Discuss the difference between probabilistic $l$-diversity and entropy $l$-diversity. Why does
				one of them appear more "diverse" than the other?}
	
	\section{Manual}
		\subsection{Running the code}
		All source code is self contained within a singular file called "e1.py".\\
		Running the application can be achieved using \textit{python3 e1.py}.\\
		This runs all the tasks and their respective outputs, Listing \ref{lst:python-output} shows the output of running the application normally.\\
		There are also python tests defined within the application. These can be run using \textit{pytest -v e1.py}.\\
		The flag "-v" is not mandatory, but improves clarity on what actually happened.\\
		Running the tests with the aforementioned command results in an output similar to Listing \ref{lst:pytest-output}.
		
		\subsection{Test Cases}
	
		\subsubsection{Task 1 Tests}
		\subsubsection{Task 2 Tests}
		\subsubsection{Task 3 Tests}
		\subsubsection{Task 4 Tests}
		\subsubsection{Task 5 Tests}
		\subsubsection{Task 6 Tests}
		\subsubsection{Task 7 Tests}
	
\end{document}